% !TeX root = main.tex
% !TeX root = main.tex
\appendix
\chapter{Requirement Specifications}
\label{app:requirements}

This appendix contains the complete list of Functional (FR) and Non-Functional (NFR) requirements defined for the Universal Speech Translation Platform.

\section{Functional Requirements}

\begin{longtable}{|l|p{10cm}|l|}
\hline
\textbf{ID} & \textbf{Description} & \textbf{Priority} \\ \hline
\endfirsthead
\hline
\textbf{ID} & \textbf{Description} & \textbf{Priority} \\ \hline
\endhead
\hline
\endfoot

FR-001 & \textbf{Session Management:} The system must support creating, starting, and stopping a session. & MUST \\ \hline
FR-002 & \textbf{Language Support:} The system must support dynamic language selection per session. MVP minimum: {da, en, de, fr, es} to English. & MUST \\ \hline
FR-003 & \textbf{Language Switching:} The system should allow changing language pairs during an active session. & SHOULD \\ \hline
FR-004 & \textbf{Microphone Stream:} The system must support live audio ingestion from a microphone via a gateway. & MUST \\ \hline
FR-005 & \textbf{Audio Upload:} The system must support uploading audio chunks as an alternative to streaming. & MUST \\ \hline
FR-006 & \textbf{Segment \& Streaming:} The system must process both discrete segments and continuous streams. & MUST \\ \hline
FR-007 & \textbf{Partial ASR Results:} ASR must produce intermediate (partial) transcriptions before the final result. & MUST \\ \hline
FR-008 & \textbf{Streaming Translation:} The translation component must update translations continuously as text evolves. & MUST \\ \hline
FR-009 & \textbf{Streaming TTS:} The system must generate and stream synthesized speech for translated text. & MUST \\ \hline
FR-010 & \textbf{Artifact Storage:} The system must store artifacts (audio, transcripts) in object storage (MinIO) for auditability. & MUST \\ \hline
FR-011 & \textbf{Retention Policy:} The system must support automatic deletion of artifacts after a configurable retention period. & MUST \\ \hline
FR-012 & \textbf{Event-Driven Orchestration:} Inter-service communication must be asynchronous via events. The workflow must be coordinated by a Pipeline Orchestrator to ensure ordering and state tracking. & MUST \\ \hline
FR-013 & \textbf{Schema Governance:} Events must be governed by schemas (e.g., Avro) with strict compatibility checks. & MUST \\ \hline
FR-014 & \textbf{Correlation:} A correlation ID must be propagated through all stages to trace user requests. & MUST \\ \hline
FR-015 & \textbf{Retry \& Degradation:} Services must retry on temporary failures and degrade gracefully on permanent ones. & MUST \\ \hline
FR-016 & \textbf{Durability:} The system must resume processing after a restart without losing queued events. & MUST \\ \hline
FR-017 & \textbf{Deduplication:} The client/UI should handle potential duplicate events (at-least-once delivery). & SHOULD \\ \hline
FR-018 & \textbf{Docker Compose:} The system must be deployable locally via Docker Compose for development. & MUST \\ \hline
FR-019 & \textbf{Docker Swarm:} The system must be deployable to a Docker Swarm cluster. & MUST \\ \hline
FR-020 & \textbf{Horizontal Scaling:} Each microservice must be capable of independent horizontal scaling. & MUST \\ \hline
FR-021 & \textbf{Ray Compute:} Compute-heavy stages (ASR, TTS) should utilize Ray Actors for warm-loaded execution. & SHOULD \\ \hline
FR-022 & \textbf{Infrastructure as Code:} Infrastructure should be provisioned via Pulumi (Post-MVP). & SHOULD \\ \hline
FR-023 & \textbf{Architectural Guardrails:} The system should implement runtime guardrails (e.g., filtering empty transcriptions, hallucination mitigation) to prevent invalid AI outputs from reaching the user. & SHOULD \\ \hline
FR-024 & \textbf{VAD as Gatekeeper:} The system must utilize VAD to segment human speech from silence before sending data to compute-heavy ASR components. & MUST \\ \hline
\end{longtable}

\section{Non-Functional Requirements}

\begin{longtable}{|l|p{10cm}|l|}
\hline
\textbf{ID} & \textbf{Description} & \textbf{Priority} \\ \hline
\endfirsthead
\hline
\textbf{ID} & \textbf{Description} & \textbf{Priority} \\ \hline
\endhead
\hline
\endfoot

NFR-001 & \textbf{Time-to-First-Output:} The first partial result must be delivered within 3 seconds. & MUST \\ \hline
NFR-002 & \textbf{End-to-End Latency:} Median processing time per segment must be $< 5$ seconds. & MUST \\ \hline
NFR-003 & \textbf{Scalability Test:} The system should support up to 20 concurrent sessions in testing. & SHOULD \\ \hline
NFR-004 & \textbf{Observability:} The system must provide metrics, tracing, and structured logging. & MUST \\ \hline
NFR-005 & \textbf{Privacy:} Logs must be anonymized and free of PII (raw audio/transcripts). & MUST \\ \hline
NFR-006 & \textbf{Encryption:} Traffic should be encrypted (TLS) in transit (Post-MVP Requirement). & SHOULD \\ \hline
NFR-007 & \textbf{Artifact Security:} Internal event references to artifacts must not contain bearer secrets. & MUST \\ \hline
NFR-008 & \textbf{Resource Efficiency:} Services should run on standard server hardware with GPU support. & SHOULD \\ \hline
NFR-009 & \textbf{Reproducibility:} Performance measurements must be reproducible via documented scripts. & MUST \\ \hline
NFR-010 & \textbf{Quality Evaluation:} Translation quality is evaluated via both automated metrics and human rating. & SHOULD \\ \hline
\end{longtable}

\chapter{Sequence Diagrams}
\label{app:sequences}

This appendix provides detailed sequence diagrams illustrating the asynchronous interaction between microservices.


\chapter{Data Contracts (Avro Schemas)}
\label{app:schemas}

This appendix documents the strict data contracts enforced across the microservice pipeline using Apache Avro. These schemas serve as the "Shared Contract Artifact" described in Section \ref{sec:data_modeling}, ensuring type safety and backward compatibility.

All events inherit a common envelope structure defined in \texttt{BaseEvent.avsc} to guarantee consistent tracing via \texttt{correlation\_id}.

\section{Base Envelope}
The base structure that all specific events must adhere to. This ensures that infrastructure-level concerns (tracing, timing, provenance) are decoupled from business payloads.

\begin{lstlisting}[language=json, caption={BaseEvent.avsc}, label={lst:base_event}]
{
  "type": "record",
  "name": "BaseEvent",
  "namespace": "ai.speech.shared.events",
  "doc": "Common envelope for all MVP events.",
  "fields": [
    {"name": "event_id", "type": "string", "doc": "UUID for this event"},
    {"name": "event_type", "type": "string", "doc": "Logical event type name"},
    {"name": "correlation_id", "type": "string", "doc": "UUID shared across pipeline"},
    {"name": "timestamp", "type": {"type": "long", "logicalType": "timestamp-millis"}, "doc": "Unix epoch ms"},
    {"name": "source_service", "type": "string", "doc": "Originating service name"}
  ]
}
\end{lstlisting}

\section{Audio Ingestion}
The \texttt{AudioInputEvent} is emitted by the Ingress Gateway. Note the \texttt{audio\_bytes} field for small payloads and \texttt{audio\_uri} for the "Claim Check" pattern if payloads exceed Kafka limits.

\begin{lstlisting}[language=json, caption={AudioInputEvent.avsc}, label={lst:audio_input}]
{
  "type": "record",
  "name": "AudioInputEvent",
  "namespace": "ai.speech.shared.events",
  "fields": [
    {"name": "event_id", "type": "string"},
    {"name": "event_type", "type": "string"},
    {"name": "correlation_id", "type": "string"},
    {"name": "timestamp", "type": {"type": "long", "logicalType": "timestamp-millis"}},
    {"name": "source_service", "type": "string"},
    {
      "name": "payload",
      "type": {
        "type": "record",
        "name": "AudioInputPayload",
        "fields": [
          {"name": "audio_bytes", "type": ["null", "bytes"], "default": null},
          {"name": "audio_uri", "type": ["null", "string"], "default": null},
          {"name": "audio_format", "type": "string", "doc": "e.g., wav"},
          {"name": "sample_rate_hz", "type": "int"},
          {"name": "language", "type": ["null", "string"], "default": null, "aliases": ["language_hint"]},
          {"name": "speaker_id", "type": ["null", "string"], "default": null}
        ]
      }
    }
  ]
}
\end{lstlisting}

\section{Speech Segmentation (VAD)}
The \texttt{SpeechSegmentEvent} represents a cleaned, verified chunk of speech ready for ASR. It includes timing information (\texttt{start\_ms}, \texttt{end\_ms}) relative to the session start.

\begin{lstlisting}[language=json, caption={SpeechSegmentEvent.avsc}, label={lst:speech_segment}]
{
  "type": "record",
  "name": "SpeechSegmentEvent",
  "namespace": "ai.speech.shared.events",
  "fields": [
    {"name": "event_id", "type": "string"},
    {"name": "correlation_id", "type": "string"},
    {
      "name": "payload",
      "type": {
        "type": "record",
        "name": "SpeechSegmentPayload",
        "fields": [
          {"name": "segment_id", "type": "string"},
          {"name": "segment_index", "type": "int"},
          {"name": "start_ms", "type": "long"},
          {"name": "end_ms", "type": "long"},
          {"name": "audio_bytes", "type": ["null", "bytes"], "default": null},
          {"name": "segment_uri", "type": ["null", "string"], "default": null},
          {"name": "sample_rate_hz", "type": "int"}
        ]
      }
    }
  ]
}
\end{lstlisting}

\section{Translation Output}
The \texttt{TextTranslatedEvent} carries the core business value: the translated text. It includes source and target language codes to ensure downstream TTS picks the correct voice.

\begin{lstlisting}[language=json, caption={TextTranslatedEvent.avsc}, label={lst:text_translated}]
{
  "type": "record",
  "name": "TextTranslatedEvent",
  "namespace": "ai.speech.shared.events",
  "fields": [
    {"name": "event_id", "type": "string"},
    {"name": "correlation_id", "type": "string"},
    {
      "name": "payload",
      "type": {
        "type": "record",
        "name": "TextTranslatedPayload",
        "fields": [
          {"name": "text", "type": ["null", "string"], "default": null},
          {"name": "source_language", "type": "string"},
          {"name": "target_language", "type": "string"},
          {"name": "quality_score", "type": ["null", "float"], "default": null}
        ]
      }
    }
  ]
}
\end{lstlisting}

\section{Synthesis Output (TTS)}
The \texttt{AudioSynthesisEvent} is the final output. Similar to the input event, it supports both inline bytes and URIs via the Claim Check pattern to handle large generated audio files.

\begin{lstlisting}[language=json, caption={AudioSynthesisEvent.avsc}, label={lst:audio_synthesis}]
{
  "type": "record",
  "name": "AudioSynthesisEvent",
  "namespace": "ai.speech.shared.events",
  "fields": [
    {"name": "event_id", "type": "string"},
    {"name": "correlation_id", "type": "string"},
    {
      "name": "payload",
      "type": {
        "type": "record",
        "name": "AudioSynthesisPayload",
        "fields": [
          {"name": "audio_bytes", "type": ["null", "bytes"], "default": null},
          {"name": "audio_uri", "type": ["null", "string"], "default": null},
          {"name": "duration_ms", "type": "long"},
          {"name": "sample_rate_hz", "type": "int"},
          {"name": "audio_format", "type": "string", "default": "wav"}
        ]
      }
    }
  ]
}
\end{lstlisting}